% sec06_1.tex — Section 6.1: Introduction
\section{Introduction}
\label{sec:fd-intro}

Partial differential equations are often classified.  Equations with the same classification have qualitatively similar mathematical and physical properties.  The heat equation ($\partial u/\partial t = k\,\partial^2 u/\partial x^2$) is an example of a \textbf{parabolic} partial differential equation.  Solutions usually exponentially decay in time and approach an equilibrium solution.  Information and discontinuities propagate at an infinite velocity.  The wave equation ($\partial^2 u/\partial t^2 = c^2\,\partial^2 u/\partial x^2$) typifies \textbf{hyperbolic} partial differential equations.  There are modes of vibration.  Information propagates at a finite velocity, and thus discontinuities persist.  Laplace's equation ($\partial^2 u/\partial x^2 + \partial^2 u/\partial y^2 = 0$) is an example of an \textbf{elliptic} partial differential equation.  Solutions usually satisfy maximum principles.  The terminology \emph{parabolic}, \emph{hyperbolic}, and \emph{elliptic} result from transformation properties of the conic sections (e.g., see Weinberger [1995]).

There are various methods to obtain explicit solutions of some partial differential equations of physical interest.  Except for one-dimensional wave equation, the solutions are rather complicated, involving an infinite series or an integral representation.  In many current situations, detailed numerical calculations of solutions of partial differential equations are needed.  Our previous analyses suggest computational methods (e.g., the first 100 terms of a Fourier series).  However, usually there are more efficient methods to obtain numerical results, especially if a computer is to be utilized.  In this chapter we develop finite difference methods to numerically approximate solutions of the different types of partial differential equations (i.e., parabolic, hyperbolic, and elliptic).  We will describe only the heat, wave, and Laplace's equations, but algorithms for more complicated problems (including nonlinear ones) will become apparent.
